\section{Tools} \label{sec3}
In this section, we introduce some results that will be used in the proof of the main theorem. Some of them are already proved in previous papers with H.-T. Yau, L. Erd\H{o}s, and A. Knowles, and we made slight changes in the statement to fit the notations and definitions in this paper. We also extend some of the known results.
 
Define the domain
\be \label{defS}
{\bf  S}(C) = \Big\{ z=E+i\eta\; : \; |E|\leq 5, \quad N^{-1}\varphi^C < \eta \le 10 \Big\}.
\ee

\begin{lem}[Previous results on generalized Wigner matrices] \label{lem: old}
Fix $\zeta>0$, $C_0\ge0$ and $c>0$. Let $H$ be a generalized Wigner matrix satisfying \eqref{eq:diag} with $C_0$ and the bounded support condition with $q$ and $c$ (Definition \ref{cond:bdd}). There exists a constant $C>0$, depending only on $(\zeta,C_0,c)$, such that, if $q\ge\varphi^C$, then the following properties hold with $\zeta$-high probability, where the constants of Definition \ref{def: hprob} depend only on $(\zeta,C_0,c)$ and not on $q$:

\begin{enumerate}
\item {\bf Local semicircle law} (Theorem 2.8 in \cite{EKYY1}): \eqref{scm} and \eqref{Gij estimate} hold simultaneously for all $z\in{\bf S}(C)$ and all $i,j$:
\begin{equation} \label{scm}
\bigcap_{z\in {\bf  S }(C) } \Big\{ |m(z)-m_{sc}(z)|
\leq \varphi^C
\left(
\min \left\{
\frac{1}{\sqrt{\kappa +\eta}}\frac{1}{q^2}, \frac{1}{q}
\right\} + \frac{1}{N \eta}
\right)
\Big\}
\end{equation}
 \begin{align}\label{Gij estimate}
\bigcap_{z\in {\bf  S }(C)  }
 \left\{
 \max_{ij} |G_{ij}(z)-\delta_{ij}m_{sc}(z)| \le
\varphi^C
\left(
\frac{1}{q} + \sqrt{\frac{\im m_{sc}(z)}{N \eta}} + \frac{1}{N \eta}
\right)
\right\}
\end{align}

\item {\bf Bound on $\|H\|$} (Lemma 4.4 in \cite{EKYY1}):

\be\label{sjsy}
\|H\|\leq 2+\varphi^C(q^{-2}+N^{-2/3})
\ee

\item {\bf Delocalization} (Remark 2.18 in \cite{EKYY1}): for every eigenvector ${\bf u}\in\mathbb R^N$ of $H$ with $\|{\bf u}\|_2=1$,

\be\label{deloc}
\max_{i} \left|{\bf u}(i)\right|^2\leq \frac{\varphi^C}{N}
\ee

\end{enumerate}

Furthermore, let $\phi>1/3$ be a constant and assume $q\ge N^{\phi}$. Then the following properties hold, where the constants of Definition \ref{def: hprob} depend only on $(\zeta,C_0,c,\phi)$, and in part 5 also on $a$, and not on $q$:

\begin{enumerate} \setcounter{enumi}{3}
\item {\bf Rigidity of the eigenvalues} (Theorem 2.13 and Remarks 2.14, 2.15 in \cite{EKYY1}): there is $C>0$, depending only on $(\zeta,C_0,c,\phi)$, such that, with $\zeta$-high probability,
\be
\bigcap_{j=1}^{N}\bigg\{|\lambda_j-\gamma_j|
{ \,\leq\,} \varphi^C \left( \Big [ \min \big ( \, j ,  N-j+1 \,  \big) \Big  ]^{-1/3} N^{-2/3} +q^{-2}\right)\bigg\} .
\ee

\item {\bf Bound on $G_{ij}$ out of spectrum} (Equations (3.32), (3.58) and (4.36)-(4.46) in \cite{EKYY1}): there is a constant $K>0$, depending only on $(\zeta,C_0,c,\phi)$, such that for every $a>12\zeta$ there is a constant $C>0$, depending on $a$ (and on $(\zeta,C_0,c,\phi)$), such that, with $\zeta$-high probability, the bounds \eqref{3838} hold simultaneously for all $E$ with
\be\label{3737}
2+\varphi^CN^{-2/3}\leq E\leq 3, \quad \eta=\varphi^{K+a} N^{-1}\kappa^{-1/2},
\ee
where $z=E+i\eta$:
\be\label{3838}
|m(z)-m_{sc}(z)|\leq \frac{1}{\varphi^aN\eta }, \quad \im m(z)\leq \frac{1}{\varphi^aN\eta },  \quad \max_{i\ne j}|G_{ij}(z)|\leq \frac{1}{\varphi^{a/3} N \eta}
\ee

\end{enumerate}

\end{lem}

\begin{proof}[Proof of Lemma \ref{lem: old}] Throughout the proof, constants depend only on $(\zeta,C_0,c)$, and in parts 4 and 5 also on $\phi$ and $a$; they do not depend on $q$, and $N\ge N_0$ with $N_0$ depending on the same parameters. The dependence of $N_0$ on $c$ is needed, since $e^{-N^c}$ must be small. By Definition \ref{cond:bdd}, $q\le\sqrt2\,N^{1/2}$. We write $\Lambda_d:=\max_i|G_{ii}-m_{sc}|$, $\Lambda_o:=\max_{i\ne j}|G_{ij}|$ and $\Lambda:=|m-m_{sc}|$, as in \cite{EKYY1}.

{\it Step 1: reduction to entries bounded almost surely.} Let $\chi_{ij}:=\mathbf 1(|h_{ij}|\le q^{-1})$ and $\mu_{ij}:=\E h_{ij}\chi_{ij}$. Set $\hat h_{ii}:=h_{ii}\chi_{ii}-\mu_{ii}$ and, for $i\ne j$, $\hat h_{ij}:=s_{ij}(h_{ij}\chi_{ij}-\mu_{ij})$, where $s_{ij}>0$ is chosen so that $\E\hat h_{ij}^2=N^{-1}$. Let $\hat H=(\hat h_{ij})$, with Green function $\hat G$. Since $\E h_{ij}=0$,
$$
|\mu_{ij}|=\big|\E h_{ij}(1-\chi_{ij})\big|\le\big(\E h_{ij}^2(1-\chi_{ij})\big)^{1/2}\,\p(|h_{ij}|>q^{-1})^{1/2}.
$$
For $i\ne j$, \eqref{eq:cond_bdd} and (2.4$'$) give $|\mu_{ij}|\le e^{-N^c}$, and $\E(h_{ij}\chi_{ij}-\mu_{ij})^2=N^{-1}-\E h_{ij}^2(1-\chi_{ij})-\mu_{ij}^2\in[N^{-1}-2e^{-N^c},N^{-1}]$, so $s_{ij}=1+O(Ne^{-N^c})$. For $i=j$, \eqref{eq:diag} and \eqref{eq:cond_bdd} give $|\mu_{ii}|\le((1+C_0)N^{-1})^{1/2}e^{-N^c/2}$. Hence:
\begin{itemize}
\item $\hat H$ is a generalized Wigner matrix satisfying \eqref{eq:diag} with the constant $\max(C_0,1)$. Indeed, the off-diagonal variances are $N^{-1}$, and $\E\hat h_{ii}^2\in[0,(1+C_0)N^{-1}]$. No tail-variance bound is assumed for the diagonal entries, so $\E\hat h_{ii}^2$ may be much smaller than $N^{-1}$; this is why the constant $\max(C_0,1)$ appears.
\item $|\hat h_{ij}|\le2q^{-1}$ almost surely, since $q^{-1}\ge(2N)^{-1/2}$ is much larger than $e^{-N^c/4}$. So $\hat H$ satisfies the bounded support condition with $q/2$ almost surely.
\item On the event $\Omega_T:=\{|h_{ij}|\le q^{-1}\ \text{for all}\ i,j\}$, whose probability is at least $1-2N^2e^{-N^c}$, we have $|h_{ij}-\hat h_{ij}|\le e^{-N^c/4}$ for all $i,j$.
\end{itemize}
On $\Omega_T$, $\|H-\hat H\|\le Ne^{-N^c/4}$. Hence the eigenvalues of $H$ and $\hat H$ differ by at most $Ne^{-N^c/4}$, and, by the resolvent identity, $|G_{ij}(z)-\hat G_{ij}(z)|\le\|H-\hat H\|\eta^{-2}\le N^3e^{-N^c/4}$ for $\eta\ge N^{-1}$; the same bound holds for $|m-\hat m|$. Every right side in Lemma \ref{lem: old} is at least $\varphi^{-O(1)}N^{-1}$ at the points $z$ considered. So these differences are absorbed by a factor $\varphi^{-1}$ of room in the corresponding bound for $\hat H$, and $2N^2e^{-N^c}\le e^{-\varphi^\zeta}$, since $\varphi^\zeta=N^{o(1)}$. The conditions $q\ge\varphi^C$ and $q\ge N^\phi$ with $\phi>1/3$ are used only up to constant factors, so the change from $q$ to $q/2$ changes only the constants. Therefore it suffices to prove parts 1, 2 and 5 when $|h_{ij}|\le q^{-1}$ almost surely and \eqref{eq:diag} holds with $\max(C_0,1)$. Parts 3 and 4 follow from parts 1 and 2 (Steps 3 and 4).

{\it Step 2: parts 1 and 2 for bounded entries.} For $C_0=0$, \eqref{scm}, \eqref{Gij estimate} and \eqref{sjsy} are Theorem 2.8 and Lemma 4.4 of \cite{EKYY1}, with the parameter $\xi$ of \cite{EKYY1} chosen as $\xi=\zeta\log\log N$, so that $(\log N)^{\xi}=\varphi^\zeta$. The probability of the exceptional event is at most $e^{-c'\varphi^\zeta}$ for some $c'>0$, uniformly in $q\ge\varphi^C$. Three points of the argument of \cite{EKYY1} are used in a form not stated there.
\begin{itemize}
\item {\it Diagonal variances.} The proofs of \cite{EKYY1} hold, with constants depending on $C_0$, under $|\E h_{ii}^2-N^{-1}|\le C_0N^{-1}$ in place of $\E h_{ii}^2=N^{-1}$. No rescaling of $G$, $\lambda$ or ${\bf u}$ by a factor $1+O(N^{-1})$ is needed, since the case of a general constant $C_0$ in \eqref{eq:diag} is treated directly. (In Section \ref{sec4} the entries of $H^S$ have variances $(1-O(N^{-1+2\epsilon}))N^{-1}$, so the factor relating them to variance $N^{-1}$ is $1+O(N^{-1+2\epsilon})$, not $1+O(N^{-1})$.)
\item {\it Large $\eta$.} The domain ${\bf S}(C)$ contains $3<\eta\le10$, which the domain of \cite{EKYY1} does not. The initial estimate of \cite{EKYY1}, Lemma 3.16, extends to $2\le\eta\le10$: there are constants $K_1,C_1,c_1>0$ such that, if $q\ge\varphi^{K_1}$, then for each $z$ with $|E|\le5$ and $2\le\eta\le10$,
$$
\p\Big(\Lambda_d(z)+\Lambda_o(z)> C_1\big(\varphi^\zeta q^{-1}+\varphi^{2\zeta}N^{-1/2}\big)\Big)\le e^{-c_1\varphi^\zeta}.
$$
The continuity argument of \cite{EKYY1} then starts at $\eta=10$.
\item {\it All $z$ at once.} The bounds \eqref{scm} and \eqref{Gij estimate} hold at each point of a lattice in ${\bf S}(C)$ with $N^{O(1)}$ points, with probability at least $1-e^{-c'\varphi^\zeta}$ each. By a union bound they hold on the whole lattice, and they extend to all $z\in{\bf S}(C)$ by Lipschitz continuity in $z$ (the derivatives of $G_{ij}$, $m$ and $m_{sc}$ are at most $\eta^{-2}\le N^2$), using $q^2\le 2N$. This gives the intersection over $z\in{\bf S}(C)$ and all $i,j$.
\end{itemize}

{\it Step 3: part 3.} Let $C_1$ and $C_2$ be the constants of parts 1 and 2, and let $C\ge\max(2C_1+3,C_2)$. Let $H{\bf u}=\mu{\bf u}$ with $\|{\bf u}\|_2=1$. By \eqref{sjsy}, $|\mu|\le3$. Let $z=\mu+i\eta$ with $\eta=\varphi^{2C_1+2}N^{-1}$, so that $z\in{\bf S}(C_1)$. Let ${\bf u}_\beta$, $\beta=1,\dots,N$, be an orthonormal eigenbasis of $H$. Then ${\bf u}$ is a combination of the ${\bf u}_\beta$ with $\lambda_\beta=\mu$, so $|{\bf u}(i)|^2\le\sum_{\beta:\lambda_\beta=\mu}|{\bf u}_\beta(i)|^2$ by the Cauchy--Schwarz inequality. Hence
$$
|{\bf u}(i)|^2\le\sum_\beta\frac{\eta^2|{\bf u}_\beta(i)|^2}{(\lambda_\beta-\mu)^2+\eta^2}=\eta\,\im G_{ii}(z).
$$
By \eqref{Gij estimate} with $q\ge\varphi^{C_1}$ and $N\eta=\varphi^{2C_1+2}$, together with $|m_{sc}|\le1$, we get $\im G_{ii}(z)\le4\le\varphi$. Therefore $|{\bf u}(i)|^2\le\varphi^{2C_1+3}N^{-1}\le\varphi^CN^{-1}$. This holds for every unit eigenvector, also when eigenvalues are repeated.

{\it Step 4: part 4.} Let $q\ge N^\phi$ with $\phi>1/3$, so that $q\ge\varphi^C$ for the constants $C$ of parts 1 and 2. For $1\le j\le N-1$, (3.6) follows from \eqref{scm} and \eqref{sjsy} as in Theorem 2.13 and Remarks 2.14, 2.15 of \cite{EKYY1}; this step is deterministic. Theorem 2.13 of \cite{EKYY1} does not cover $j=N$, which we treat separately. By \eqref{def:gamma}, $\gamma_N=2$ and $2-\gamma_{N-1}\le C'N^{-2/3}$, since $\varrho_{sc}(x)\ge c''\sqrt{2-x}$ for $0\le x\le2$. Then \eqref{sjsy} gives $\lambda_N-2\le\varphi^C(q^{-2}+N^{-2/3})$. Moreover, $\lambda_N\ge\lambda_{N-1}$, and (3.6) at $j=N-1$ gives
$$
\lambda_N\ge\gamma_{N-1}-\varphi^C(N^{-2/3}+q^{-2})\ge2-\varphi^{C+1}(N^{-2/3}+q^{-2}).
$$
This is (3.6) at $j=N$, with $C+1$ in place of $C$.

{\it Step 5: part 5.} Let $a>12\zeta$. By Step 1 we may assume $|h_{ij}|\le q^{-1}$ almost surely.

{\it The first two bounds of \eqref{3838}.} In (4.36)--(4.46) of \cite{EKYY1} it is proved that, for $E$ in the window
$$
2+\varphi^{C'}(q^{-2}+N^{-2/3})\le E\le3
$$
of (4.36) of \cite{EKYY1}, one has $\Lambda\le((\log N)N\eta)^{-1}$ and $\im m\le((\log N)N\eta)^{-1}$ (see (4.45) and (4.46) of \cite{EKYY1}); here $\alpha\sim\sqrt\kappa$ in the notation of \cite{EKYY1}. Run at level $\zeta$, the same argument gives the following. There is $K_0>0$, depending only on $(\zeta,C_0)$, such that for every $a'>0$ there is $C'$ for which, with $\zeta$-high probability,
$$
\Lambda(z)\le\frac{1}{\varphi^{a'}N\eta},\qquad \im m(z)\le\frac{1}{\varphi^{a'}N\eta},\qquad \eta=\varphi^{K_0+a'}N^{-1}\kappa^{-1/2},
$$
for all $E$ in this window. To see this, replace $\varphi^3(\log N)^2$ in (4.37) of \cite{EKYY1} by $\varphi^{K_0+2a'}$, and replace $\varphi^3\log N$ in (4.38), (4.42) and below (4.44) of \cite{EKYY1} by $\varphi^{K_0+a'}$. The offset $K_0$ collects the constants of \cite{EKYY1} that enter at level $\zeta$. In the original statement the offset was $3$, obtained by setting $C_1=1$ in (2.15) of \cite{EKYY1}; this conflates two different constants $C_1$ of \cite{EKYY1}, so the offset is a constant $K_0$ instead.

The bound holds at each fixed $E$ with $\zeta$-high probability. It holds simultaneously for all $E$ in the window by a union bound over a grid of spacing $N^{-10}$, because $m(E+i\eta(E))$, $m_{sc}(E+i\eta(E))$ and the right sides have derivatives in $E$ bounded by $N^{O(1)}$ on the window (there $\eta(E)\ge N^{-1}$ and $|\eta'(E)|\le1$).

Since $q\ge N^\phi$ with $\phi>1/3$, we have $q^{-2}\le N^{-2/3}$. Hence the window \eqref{3737} with $C\ge C'+1$ lies inside the window above. Set $K:=K_0+K_2+1$, where $K_2$ is the constant of the next paragraph, and apply the above with $a'=a+K-K_0$. Then $\eta$ is the $\eta$ of \eqref{3737}, and the bounds hold with the extra factor $\varphi^{-(K-K_0)}\le\varphi^{-1}$, which leaves room for Step 1. This proves the first two parts of \eqref{3838}.

{\it The third bound of \eqref{3838}.} On the window \eqref{3737} we have $\varphi^{K}N^{-1}<\eta\le1$, so $z\in{\bf S}(K_2)$ once $K\ge K_2$. Using (3.32) and (3.58) of \cite{EKYY1} at level $\zeta$, there are $K_2,C_o>0$, depending only on $(\zeta,C_0)$, such that, if $q\ge\varphi^{K_2}$, then with $\zeta$-high probability
$$
\Lambda_o(z)\le C_o\big(q^{-1}+\varphi^{2\zeta}\Psi(z)\big),\qquad \Psi:=\sqrt{\frac{\Lambda+\im m_{sc}}{N\eta}},
$$
where $\Psi$ is the control parameter of (3.26) of \cite{EKYY1}. On the event of the first two parts, $\im m_{sc}\le\im m+\Lambda\le2(\varphi^aN\eta)^{-1}$, so $\Psi\le\sqrt3\,\varphi^{-a/2}(N\eta)^{-1}$. Renaming $\sqrt3C_o$ as $C_o$, we get, with $\zeta$-high probability,
\be\label{tyzz}
\max_{i\ne j} |G_{ij}| \leq C_o\Big(\frac1q+\frac{\varphi^{2\zeta}}{\varphi^{a/2} N \eta}\Big).
\ee
For the second term, $C_o\varphi^{2\zeta-a/2}\le\frac13\varphi^{-a/3}$ holds for $N\ge N_0$ if $a/6-2\zeta>0$, that is, $a>12\zeta$. At $\zeta=1$ this is the condition $a>12$. For the first term, \eqref{3737} gives $\kappa=E-2\ge\varphi^CN^{-2/3}$, hence
$$
N\eta=\varphi^{K+a}\kappa^{-1/2}\le\varphi^{K+a-C/2}N^{1/3},
$$
and so $\varphi^{a/3}N\eta\,C_oq^{-1}\le C_o\varphi^{K+4a/3}N^{1/3-\phi}\le\frac13$ for $N\ge N_0$, since $\phi>1/3$. Hence $\max_{i\ne j}|G_{ij}|\le\frac23(\varphi^{a/3}N\eta)^{-1}$ at each $E$ with $\zeta$-high probability. The grid argument above makes this simultaneous in $E$, and the remaining factor $\frac23$ leaves room for Step 1. This gives the third part of \eqref{3838} and completes the proof.

\end{proof}

\begin{rem} \label{prozeta}
Lemma \ref{lem: old} holds for every fixed $\zeta>0$, with a constant $C=C_\zeta$; it is stated in this form. In this paper we use $\zeta\in\{1,3,4\}$; the case $\zeta=4$ is used for \eqref{3838T}.
\end{rem}
